% ---
% id: EX-CH01-007
% titre: PGDC et PPMC
% chapitre_id: 1
% chapitre_nom: Nombres entiers & divisibilité
% annee_scolaire: 9H
% difficulte: 2
% tags: [divisibilite, pgdc, ppmc]
% auteur: Bussy D.
% ---

\begin{exercice}{EX-CH01-007}{PGDC et PPMC}

Recherche le PGDC et le PPMC des nombres suivants :

\bigskip
\textbf{A.}

\medskip
\sol{
\def\arraystretch{1.7}
\begin{tabularx}{\linewidth}{>{\centering\arraybackslash}XXXXp{4cm}>{\centering\arraybackslash}XXXXp{4cm}}
\textbf{L.} & 34 & et & 26    && \textbf{M.} & 130 & et & 260 & \\
\textbf{O.} & 25 & et & 45    && \textbf{A.} & 108 & et & 225 & \\
\textbf{V.} & 16 & et & 18    && \textbf{T.} & 375 & et & 2070 & \\
\textbf{E.} & 120 & et & 300  && \textbf{H.} & 12 & et & 8 & \\
\end{tabularx}}
{\footnotesize
\medskip
\def\arraystretch{1.7}
\begin{tabularx}{\linewidth}{>{\centering\arraybackslash}lp{6.3cm}XX}
\textbf{L.} & $34 = 2 \times 17$ et $26 = 2 \times 13$ &PGDC(34,26) = 2& PPMC(34,26) = 442 \\
\textbf{O.} &$25 = 5^2$ et $45 = 3^2 \times 5$&PGDC(25,45) = 5&PPMC(25,45) = 225 \\
\textbf{V.} &$16 = 2^4$ et $18 = 2 \times 3^2$&PGDC(16,18) = 2&PPMC(16,18) = 144 \\
\textbf{E.} &$120 = 2^3 \times 3 \times 5$ et $300 = 2^2 \times 3 \times 5^2$&PGDC(120,300) = 60&PPMC(120,300) = 600 \\
\end{tabularx}

\medskip
\begin{tabularx}{\linewidth}{>{\centering\arraybackslash}lp{6.3cm}XX}
\textbf{M.} &$130 = 2 \times 5 \times 13$ et $260 = 2^2 \times 5 \times 13$&PGDC(130,260) = 130&PPMC(130,260) = 260 \\
\textbf{A.} &$108 = 2^2 \times 3^3$ et $225 = 3^2 \times 5^2$&PGDC(108,225) = 9&PPMC(108,225) = 2700 \\
\textbf{T.} &$375 = 3 \times 5^3$ et $2070 = 2 \times 3^2 \times 5 \times 23$&PGDC(375,2070) = 15&PPMC(375,2070) = 51750 \\
\textbf{H.} &$12 = 2^2 \times 3$ et $8 = 2^3$&PGDC(12,8) = 4&PPMC(12,8) = 24 \\
\end{tabularx}
\def\arraystretch{1}}

\vspace{20pt}
\textbf{B.}

\medskip
\sol{
\def\arraystretch{1.7}
\begin{tabularx}{\linewidth}{>{\centering\arraybackslash}XXXXp{4cm}>{\centering\arraybackslash}XXXXXp{3cm}}
\textbf{L.} & 382 & et & 675&& \textbf{M.} & 180&  et &378&& \\
\textbf{I.} & 528 & et  & 204&& \textbf{A.} & 384 & et &132&& \\
\textbf{E.} & 25 & et  & 16&& \textbf{T.} & 45& et &90&& \\
\textbf{B.} & 54 & et & 36&& \textbf{H.} & 384,& 126& et &132& \\
\end{tabularx}}
{\footnotesize
\medskip
\def\arraystretch{1.7}
\begin{tabularx}{\linewidth}{>{\centering\arraybackslash}lp{6.3cm}XX}
\textbf{L.} &$382 = 2 \times 191$ et $675 = 3^3 \times 5^2$&PGDC(382,675) = 1&PPMC(382,675) = 257850 \\
\textbf{I.} &$528 = 2^4 \times 3 \times 11$ et $204 = 2^2 \times 3 \times 17$&PGDC(528,204) = 12&PPMC(528,204) = 8976 \\
\textbf{E.} &$25 = 5^2$ et $16 = 2^4$&PGDC(25,16) = 1& PPMC(25,16) = 400\\
\textbf{B.} &$54 = 2 \times 3^3$  et $36 = 2^2 \times 3^2$&PGDC(54,36) = 18& PPMC(54,36) = 108\\
\end{tabularx}

\medskip
\def\arraystretch{1.7}
\begin{tabularx}{\linewidth}{>{\centering\arraybackslash}lp{6.3cm}XX}
\textbf{M.} &$180 = 2^2 \times 3^2 \times 5$ et $378 = 2 \times 3^3 \times 7$&PGDC(180,378) = 18&PPMC(180,378) = 3780 \\
\textbf{A.} &$384 = 2^7 \times 3$et $132 = 2^2 \times 3 \times 11$&PGDC(384,132) = 12&PPMC(384,132) = 4224 \\
\textbf{T.} &$45 = 3^2 \times 5$ et $90 = 2 \times 3^2 \times 5$&PGDC(45,90) = 45&PPMC(45,90) = 90 \\
\textbf{H.} &$384 = 2^7 \times 3$ et $126 = 2 \times 3^2 \times 7$ et $132 = 2^2 \times 3 \times 11$&PGDC(384,126,132) = 6& PPMC(384,126,132) = 21504\\
\end{tabularx}
\def\arraystretch{1}}
\end{exercice}